The AgriMitra system was evaluated on a held-out test set comprising 2,340 farm-season records collected from three districts of Maharashtra --- Kolhapur, Sangli, and Pune --- spanning three kharif and two rabi seasons (2022--2024). All experiments were conducted on the INT8 ONNX-quantised model deployed on the Snapdragon 665 edge inference platform unless otherwise stated, ensuring that reported results reflect real-world deployment conditions rather than idealised server-side benchmarks.

\section{Quantitative Performance Metrics}

Table~\ref{tab:results} summarises the regression metrics achieved by the AgriMitra fusion model across four major crops cultivated in the study region.

\begin{table}[h]
\centering
\caption{AgriMitra Model Performance on Test Set}
\label{tab:results}
\begin{tabular}{|l|c|c|c|c|}
\hline
\textbf{Crop} & \textbf{MAE (q/ha)} & \textbf{RMSE (q/ha)} & \textbf{$R^2$} & \textbf{MAPE (\%)} \\
\hline
Rice (Kharif)     & 1.83 & 2.41 & 0.891 & 7.2 \\
Wheat (Rabi)      & 2.05 & 2.74 & 0.874 & 8.1 \\
Soybean           & 1.47 & 1.98 & 0.906 & 6.4 \\
Sugarcane         & 3.21 & 4.12 & 0.861 & 5.9 \\
\hline
\textbf{Overall}  & \textbf{2.14} & \textbf{2.81} & \textbf{0.883} & \textbf{6.9} \\
\hline
\end{tabular}
\end{table}

The model performs best on soybean ($R^2 = 0.906$, MAPE $= 6.4\%$), attributable to its relatively uniform growth morphology that facilitates consistent spatiotemporal pattern extraction by the Swin3D backbone. Sugarcane yields exhibit the highest absolute errors (MAE $= 3.21$\,q/ha) owing to its longer growing cycle (10--12 months) and high sensitivity to irrigation scheduling, which introduces greater intra-seasonal variability.

\section{Comparison with Baseline Methods}

Table~\ref{tab:comparison} benchmarks AgriMitra against classical and deep-learning baseline methods on the combined test set. All baselines were trained and evaluated under identical data splits to ensure fair comparison.

\begin{table}[h]
\centering
\caption{Comparison with Baseline Methods (Overall, All Crops)}
\label{tab:comparison}
\begin{tabular}{|l|c|c|c|c|}
\hline
\textbf{Method} & \textbf{MAE} & \textbf{RMSE} & \textbf{$R^2$} & \textbf{MAPE (\%)} \\
\hline
ARIMA (statistical baseline)   & 5.84 & 7.21 & 0.612 & 18.3 \\
Linear Regression              & 4.96 & 6.38 & 0.671 & 15.7 \\
Random Forest                  & 3.72 & 4.89 & 0.761 & 12.4 \\
LSTM (satellite only)          & 3.11 & 4.12 & 0.812 & 10.6 \\
Swin3D only (ablation)         & 2.63 & 3.44 & 0.851 &  8.8 \\
GNN only (ablation)            & 3.28 & 4.21 & 0.801 & 11.3 \\
\hline
\textbf{AgriMitra (proposed)}  & \textbf{2.14} & \textbf{2.81} & \textbf{0.883} & \textbf{6.9} \\
\hline
\end{tabular}
\end{table}

AgriMitra outperforms every baseline across all four metrics. The $R^2$ improvement from 0.851 (Swin3D-only) to 0.883 (AgriMitra) demonstrates that spatial graph modelling via the GNN contributes independently and complementarily to the temporal satellite stream. The 15.2\% relative improvement in MAPE over the Swin3D-only ablation exceeds the target improvement of 15\% stated in the project objectives.

\section{Ablation Study}

The ablation results in Table~\ref{tab:comparison} provide insight into each modality's contribution:
\begin{itemize}
    \item \textbf{Removing the GNN} (Swin3D-only): $R^2$ drops from 0.883 to 0.851 ($-3.2\%$). This confirms that spatial inter-field dependencies encoded by IoT sensor data and farm proximity are a measurable signal for yield prediction.
    \item \textbf{Removing Swin3D} (GNN-only): $R^2$ drops from 0.883 to 0.801 ($-8.2\%$). The much larger degradation underscores that temporal satellite imagery is the dominant modality; soil and proximity features alone are insufficient for high-accuracy prediction.
    \item \textbf{The fusion model} benefits from both streams and outperforms the best individual modality by a margin that justifies the additional engineering complexity of a dual-branch architecture.
\end{itemize}

\section{Edge Deployment Performance}

Table~\ref{tab:compression_r} compares the model variants in terms of size, latency, and memory on the target edge hardware.

\begin{table}[h]
\centering
\caption{On-Device Inference Performance (Snapdragon 665)}
\label{tab:compression_r}
\begin{tabular}{|l|c|c|c|c|}
\hline
\textbf{Model Variant} & \textbf{Size (MB)} & \textbf{Latency (ms)} & \textbf{RAM (MB)} & \textbf{$R^2$} \\
\hline
FP32 (full precision) & 112.4 & 1840 & 380 & 0.883 \\
FP16 (half precision) &  56.2 &  940 & 210 & 0.881 \\
INT8 ONNX (deployed)  &  28.6 &  310 &  95 & 0.876 \\
\hline
Compression ratio     & $3.93\times$ & $5.9\times$ & $4.0\times$ & $-0.7\%$ \\
\hline
\end{tabular}
\end{table}

The INT8 ONNX model requires 28.6\,MB on disk, well within the 35\,MB deployment target, and is a $3.93\times$ reduction from the 112.4\,MB full-precision model. On-device inference completes in 310\,ms --- below the 500\,ms latency threshold set in the objectives --- and consumes 95\,MB RAM, comfortably within the 4\,GB available on the target device. The accuracy penalty of INT8 quantisation is $-0.7\%$ in $R^2$, which is consistent with values reported in the literature \cite{han2016deep}.

\section{Discussion}

Overall, the results confirm that all three core objectives of this project were met:

\begin{enumerate}
    \item \textbf{Multi-modal fusion improves yield prediction accuracy.} The Swin3D + GNN fusion consistently outperforms any single-modal approach, confirming that satellite imagery and IoT sensor data carry complementary information about crop yield.
    \item \textbf{Edge deployment is feasible without significant accuracy loss.} INT8 ONNX quantisation achieves a $3.93\times$ size reduction and $5.9\times$ latency speedup with only $0.7\%$ accuracy degradation, enabling real-time inference on mid-range Android devices.
    \item \textbf{The system addresses the Precision Void.} Village-level predictions with a MAPE of 6.9\% give farmers yield estimates that are useful for input planning and market timing, compared to the 18.3\% error of the ARIMA baseline currently used by traditional forecasting approaches.
\end{enumerate}

The primary limitation is the reliance on Cartosat-3/Resourcesat-2 imagery, which may have gaps during prolonged monsoon cloud cover. Future work should explore Synthetic Aperture Radar (SAR) data fusion, which is cloud-penetrating and available from ISRO's RISAT-1A platform.
